\section{Solutions for Chapter~\ref{ch:arithmetic}}

\paragraph{Exercise~\ref{ex:HI-periods}.}
\phantomsection
\label{sol:ex:HI-periods}
Dividing $d_n^2=nd_n+1$ by $d_n$ gives
$d_n-d_n^{-1}=n$, hence the period difference.
Expanding $(d_n+1)(1+d_n^{-1})$ gives
$d_n+2+d_n^{-1}=\omega_1+\omega_2$; division by their
product proves the reciprocal identity.  For nonzero integer
$a$,
\[
 \sin(\pi a/\omega_1)
 =\sin(\pi a-\pi a/\omega_2)
 =(-1)^{a+1}\sin(\pi a/\omega_2).
\]
The denominator is nonzero because $\omega_2$ is irrational.
Substituting this sine ratio into
$F(a+n)/F(a)=\sin(\pi a/\omega_1)/\sin(\pi a/\omega_2)$
gives $F(a+n)=(-1)^{a+1}F(a)$.
At the origin the sine ratio tends instead to
$\omega_2/\omega_1=d_n^{-1}$.



\paragraph{Exercise~\ref{ex:HI-sign}.}
\phantomsection
\label{sol:ex:HI-sign}
For $1\leq a<n$, multiplication of the two sample formulas gives
the sign exponent
\[
 \frac{a(a+n)+(n-a)(2n-a)}2=a(a-n)+n^2.
\]
Since $n$ is odd, $a$ and $a-n$ have opposite parity, so their
product is even; $n^2$ is odd. The two initial minus signs cancel,
the remaining sign is minus, and reflection gives
$W_aW_{n-a}=-F(a)F(n-a)/d_n=-d_n^{-1}$.
The integer exponent is not itself a function of the residue class:
replacing $a$ by $a+n$ changes $a(a+n)/2$ by $n(a+n)$,
which can be odd. For example at $n=3$, representatives $2$ and
$5$ give exponents $5$ and $20$. One must state the chosen
representatives and include the accompanying gamma-shift factor
when changing them. The separately defined origin $W_0=-1$
also cannot be obtained by the generic nonzero-sample rule.



\paragraph{Exercise~\ref{ex:exact-rank}.}
\phantomsection
\label{sol:ex:exact-rank}
Take
\[
 M=\begin{pmatrix}1&1\\1&1+10^{-8}\end{pmatrix}.
\]
Its exact determinant is $10^{-8}$, while rounding to six decimal
places produces two identical rows.  After clearing denominators,
a nonzero determinant modulo a prime proves that the original
integer determinant was nonzero.  A zero residue only says that
the prime divides that determinant; the integer could still be
nonzero.  Modular rank methods therefore give decisive lower
bounds when they succeed and potentially inconclusive results
when they fail.
For an explicit failure at a prime with no denominator problem,
$\operatorname{diag}(1,7)$ has exact determinant seven but
determinant zero modulo seven. A nonzero $r\times r$ minor
modulo a permitted prime proves rank at least $r$ over $\Q$;
proving an upper bound still requires showing the larger minors
vanish exactly.



\paragraph{Exercise~\ref{ex:component-norm}.}
\phantomsection
\label{sol:ex:component-norm}
The two evaluations of $1+t$ are $2$ and $0$, so their product
is zero and the element is not invertible.  The two evaluations
of $2+t$ are $3$ and $1$, so its product is three.  It is
invertible, with inverse $(2-t)/3$.  Modulo five both conclusions
remain valid: the relevant products are zero and three.
The distinction is visible in the algebra itself:
$(1+t)(1-t)=0$ with $1-t\ne0$, whereas
$(2+t)(2-t)=4-t^2=3$. Since two and three remain invertible
modulo five, the algebra still splits into its $t=1$ and $t=-1$
components, and the second inverse reduces to $4-2t$.



\paragraph{Exercise~\ref{ex:unitarity-hierarchy}.}
\phantomsection
\label{sol:ex:unitarity-hierarchy}
A field automorphism preserves a polynomial equality by preserving
addition and multiplication.  Positivity is an order condition,
and an adjoint involves complex conjugation, neither of which is
part of that algebraic guarantee.  For example the two real
embeddings of $\Q(\sqrt{13})$ send the positive number
$\sqrt{13}$ to opposite signs.  Thus a transported algebraic
solution must have its Hermitian and positivity properties
checked separately.


